\documentclass[12pt,a4paper]{article}

\usepackage[utf8]{inputenc}
\usepackage[spanish]{babel}
\usepackage{graphicx}
\usepackage{geometry}
\geometry{a4paper, margin=2.5cm}
\usepackage{xcolor}
\usepackage{listings}
\usepackage{booktabs}
\usepackage{amsmath}
\usepackage[colorlinks=true, linkcolor=blue, urlcolor=blue]{hyperref}

\definecolor{codegreen}{rgb}{0,0.6,0}
\definecolor{codegray}{rgb}{0.5,0.5,0.5}
\definecolor{codepurple}{rgb}{0.58,0,0.82}
\definecolor{backcolour}{rgb}{0.95,0.95,0.92}

\lstdefinestyle{mystyle}{
    backgroundcolor=\color{backcolour},   
    commentstyle=\color{codegreen},
    keywordstyle=\color{magenta},
    numberstyle=\tiny\color{codegray},
    stringstyle=\color{codepurple},
    basicstyle=\ttfamily\footnotesize,
    breakatwhitespace=false,         
    breaklines=true,                 
    captionpos=b,                    
    keepspaces=true,                 
    numbers=left,                    
    numbersep=5pt,                  
    showspaces=false,                
    showstringspaces=false,
    showtabs=false,                  
    tabsize=2
}
\lstset{style=mystyle}

\title{\Huge\textbf{PROYECTO ESP32} \\ \Large On/Off Repetitivo \\ \vspace{0.5cm} \normalsize \textit{Parpadeo automático de un LED usando el microcontrolador ESP32}}
\author{Miguel Angel Delagado Acevedo - Juan José Peres - Daisy - Andres Felipe}
\date{2026 \\ \vspace{0.2cm} \small Plataforma: ESP32 Dev Module}

\begin{document}

\maketitle
\tableofcontents
\newpage

\section{Introducción}
Este documento describe de manera detallada el proyecto "On/Off Repetitivo", un circuito básico pero fundamental en el aprendizaje de la programación de microcontroladores. El proyecto consiste en hacer parpadear un LED conectado al pin GPIO 4 del ESP32, alternando entre estados de encendido y apagado con intervalos de 1 segundo.

El ESP32 es un microcontrolador de bajo costo y alto rendimiento desarrollado por Espressif Systems. Cuenta con conectividad Wi-Fi y Bluetooth integrada, lo que lo convierte en una plataforma ideal tanto para proyectos educativos como para aplicaciones IoT profesionales.

Este proyecto sirve como punto de partida para comprender conceptos clave como: configuración de pines GPIO, uso de funciones de tiempo (delay), comunicación serial para depuración, y el flujo básico de un programa Arduino (setup/loop).

\section{Objetivo del Proyecto}
Implementar un sistema de parpadeo automático de un LED utilizando el microcontrolador ESP32, demostrando el uso correcto de salidas digitales, la estructura de un programa Arduino, y la comunicación serial como herramienta de monitoreo.

\section{Materiales Necesarios}
\begin{table}[h!]
\centering
\begin{tabular}{cll}
\toprule
\textbf{Cant.} & \textbf{Componente} & \textbf{Especificación} \\
\midrule
1 & ESP32 Dev Module & Microcontrolador con Wi-Fi/Bluetooth \\
1 & LED rojo (5mm) & LED difuso, $V_f \approx 2V$, $I_f \approx 20mA$ \\
1 & Resistencia 220 $\Omega$ & 1/4W, tolerancia 5\% \\
1 & Protoboard & Placa de pruebas de 830 puntos \\
2 & Cables jumper & Macho-macho para conexiones \\
1 & Cable USB Micro-B & Para programación y alimentación \\
\bottomrule
\end{tabular}
\caption{Lista de materiales}
\end{table}

\section{Esquema de Conexión}
El siguiente esquema, realizado en Fritzing, muestra las conexiones físicas del circuito. El LED rojo está conectado al pin D4 (GPIO4) del ESP32 a través de una resistencia limitadora de corriente de 220 $\Omega$. El cátodo del LED se conecta a la línea de tierra (GND) del ESP32.

\begin{figure}[h!]
\centering
\includegraphics[width=0.9\textwidth]{../Esquemas/On_Off_Repetitivo_esq/On_Off_Repetitivo_esq.jpeg}
\caption{Esquema de conexión del circuito On/Off Repetitivo (Fritzing)}
\end{figure}

\subsection{Descripción de las Conexiones}
\begin{enumerate}
    \item Pin D4 (GPIO4) del ESP32 se conecta al ánodo (+) del LED mediante un cable jumper azul.
    \item El cátodo (-) del LED se conecta a un extremo de la resistencia de 220 $\Omega$.
    \item El otro extremo de la resistencia se conecta al bus de tierra (GND) de la protoboard.
    \item El pin GND del ESP32 se conecta al bus de tierra de la protoboard con un cable negro.
\end{enumerate}

\section{Código Fuente}
A continuación se presenta el código completo del proyecto con explicación detallada de cada sección:

\subsection{Código Completo (On\_Off\_Repetitivo.ino)}
\begin{lstlisting}[language=C++, title=On\_Off\_Repetitivo.ino]
/*
 * ====================================================
 *  Parpadeo de un LED conectado al pin D4 del ESP32
 *  Autor : Msc. Nestor Fabio Montoya Palacios
 *  Fecha : 2026
 *  Board : ESP32 Dev Module
 * ====================================================
 */

// -- Configuracion ----------------------------------
const int PIN_LED        = 4;     // Pin GPIO donde esta el LED
const int TIEMPO_ENCENDIDO  = 1000;  // Tiempo encendido  (ms)
const int TIEMPO_APAGADO    = 1000;  // Tiempo apagado    (ms)

// -- Setup -------------------------------------------
void setup() {
  Serial.begin(115200);
  Serial.println("==================================");
  Serial.println("  Parpadeo de LED con el ESP32    ");
  Serial.println("==================================");

  // Establecemos el pin del LED como salida
  pinMode(PIN_LED, OUTPUT);

  // Apagamos el LED al iniciar
  digitalWrite(PIN_LED, LOW);
  Serial.println("> Sistema listo. Iniciando parpadeo...");
}

// -- Loop --------------------------------------------
void loop() {
  // Encendemos el LED
  digitalWrite(PIN_LED, HIGH);
  Serial.println("LED: ENCENDIDO");
  delay(TIEMPO_ENCENDIDO);

  // Apagamos el LED
  digitalWrite(PIN_LED, LOW);
  Serial.println("LED: APAGADO");
  delay(TIEMPO_APAGADO);
}
\end{lstlisting}
\textit{Nota: Se han modificado los caracteres especiales y emojis del código original para garantizar la compatibilidad con el compilador de \LaTeX.}

\subsection{Análisis del Código}

\subsubsection*{Sección de Configuración (Constantes)}
Al inicio del programa se definen tres constantes que parametrizan el comportamiento del circuito:
\begin{itemize}
    \item \textbf{\texttt{PIN\_LED = 4}:} Define el pin GPIO 4 como salida para controlar el LED. En el ESP32, GPIO4 corresponde al pin físico D4.
    \item \textbf{\texttt{TIEMPO\_ENCENDIDO = 1000}:} Tiempo en milisegundos que el LED permanecerá encendido (1 segundo).
    \item \textbf{\texttt{TIEMPO\_APAGADO = 1000}:} Tiempo en milisegundos que el LED permanecerá apagado (1 segundo).
\end{itemize}

\subsubsection*{Función \texttt{setup()}}
Esta función se ejecuta una sola vez al iniciar el ESP32:
\begin{itemize}
    \item \textbf{\texttt{Serial.begin(115200)}:} Inicializa la comunicación serial a 115200 baudios para monitorear el estado del LED desde el computador.
    \item \textbf{\texttt{pinMode(PIN\_LED, OUTPUT)}:} Configura el pin GPIO4 como salida digital, permitiendo enviar señales HIGH o LOW para controlar el LED.
    \item \textbf{\texttt{digitalWrite(PIN\_LED, LOW)}:} Asegura que el LED inicie en estado apagado.
\end{itemize}

\subsubsection*{Función \texttt{loop()}}
Esta función se ejecuta en un ciclo infinito y contiene la lógica principal del parpadeo:
\begin{enumerate}
    \item \texttt{digitalWrite(PIN\_LED, HIGH)}: Enciende el LED enviando 3.3V al pin GPIO4.
    \item \texttt{Serial.println()}: Imprime el estado actual del LED en el Monitor Serial.
    \item \texttt{delay(TIEMPO\_ENCENDIDO)}: Pausa la ejecución por 1000ms (1 segundo).
    \item \texttt{digitalWrite(PIN\_LED, LOW)}: Apaga el LED enviando 0V al pin.
    \item \texttt{delay(TIEMPO\_APAGADO)}: Pausa por otro segundo antes de repetir el ciclo.
\end{enumerate}

\section{Principio de Funcionamiento}
El circuito opera bajo el principio básico de control digital de un GPIO (General Purpose Input/Output). Cuando el ESP32 escribe un valor HIGH en el pin GPIO4, este entrega aproximadamente 3.3V, lo cual permite que la corriente fluya a través del LED y la resistencia, encendiendo el LED.

La resistencia de 220 $\Omega$ cumple la función crítica de limitar la corriente que circula por el LED. Sin esta resistencia, la corriente excesiva dañaría el LED. El cálculo de la corriente es:

\begin{equation}
I = \frac{V_{gpio} - V_{led}}{R} = \frac{3.3\text{ V} - 2.0\text{ V}}{220\,\Omega} \approx 5.9\text{ mA}
\end{equation}

Este valor está dentro del rango seguro de operación tanto para el LED (máx 20mA) como para el pin del ESP32 (máx 40mA por pin).

\section{Video de Funcionamiento}
El proyecto incluye un video demostrativo (\texttt{On\_Off\_Repetitivo\_Funcionamiento.mp4}) que muestra el circuito en operación. En el video se puede observar el LED parpadeando con intervalos regulares de 1 segundo, confirmando el correcto funcionamiento del código y las conexiones del circuito.

\textit{Nota: El archivo de video se encuentra en la carpeta \texttt{Videos/On\_Off\_Repetitivo\_vid/} del proyecto.}

\section{Conclusiones}
\begin{itemize}
    \item El proyecto demuestra el uso correcto de los pines GPIO del ESP32 como salidas digitales para controlar dispositivos externos.
    \item La comunicación serial permite monitorear en tiempo real el estado del sistema, facilitando la depuración.
    \item El uso de constantes para definir pines y tiempos hace el código modular y fácil de modificar.
    \item La resistencia limitadora es esencial para proteger tanto el LED como el microcontrolador.
    \item Este proyecto constituye la base para circuitos más complejos como semáforos, indicadores de estado, y sistemas de alerta visual.
\end{itemize}

\end{document}
